\section{The Menzerath-Altman law in the verbal domain}\label{MAL}
MAL predicts a negative correlation between constituents’ average length and their total number in any given domain. MAL has been previously studied in the verbal domain using the UD annotation schema in  Czech \citep{macutek2017menzerath}, as well as  in the 2.3 version of the UD collection with 76 languages \citep{tanaka2021menzerath}.  
Our study builds on the latter while proposing a number of substantial improvements in order to be able to draw more reliable and typologically valid conclusions. First, we study MAL in the preverbal and postverbal domains sepertayly, in addition to the bilateral MAL. Second, our data extraction is more fine-grained than in Tanaka-Ishii’s study (see Section \ref{UD}). Third, we use the latest version of the UD, with 180 languages, offers a more important diversity of languages around the world. Finally, we analyse the data while taking into account language families as well as language types, that is, VO, OV or no dominant order (NDO).

%proposes substantial differences with Tanaka-Ishii’s study:
%\begin{itemize}
    %\item we filter verbal dependents the verbal constructions to retain only syntactic dependents and to eliminate coordinated conjuncts or discourse markers;
    %\item we also estimate the mean size of constituents on the right or on the left to see whether MAL is still verified on the half construction; 
    %\item we extend this research  by studying the 2.17 version of the UD collection which, with 180 languages, offers a more important diversity of languages around the world; 
    %\item we also have a discussion on the choice of the right metrics.
%\end{itemize}

